%%%%%%%%%%%%%%%%%%%%%%%%%%%%%%%%%%%%%%%%%%%%%%%%%%%%%%%%%%%%%%%%%%
% The following comments were written in Portuguese, because this 
% template applies only for Institute of Philosophy and the Humanities at University 
% of Campinas, Brazil.
%
% Este é um modelo Latex para monografias de  Mestrado e Teses de doutorado em Demografia do Instituto de Filosofia e Ciências Humanas da Universidade 
% Estadual de Campinas (UNICAMP).
%
% Este Modelo foi adaptado do Modelo criado pelo autor André Leon Sampaio Gradvohl (gradvohl@ft.unicamp.br), Dr. para a FT da UNICAMP

% Autor: Jose H C Monteiro da Silva
% Email: j117439@dac.unicamp.br
%  
% Última versão: 06/Março/2019
%
% Primeira versão (06/Março/2019): adaptação do modelo inicial desenvolvido pelo André para a FT de modo a adequá-lo às normas do IFCH.

%%%%%%%%%%%%%%%%%%%%%%%%%%%%%%%%%%%%%%%%%%%%%%%%%%%%%%%%%%%%%%%%%%
%
% Escolha: Portugues ou Ingles ou Espanhol.
% Para a versão final do texto, acrescente a palavra "Final",
% como na segunda linha abaixo da próxima.
%\documentclass[Portugues]{tese-FT}
\documentclass[Portugues,Final]{tese-IFCH}
%\documentclass[Ingles]{tese-FT}
%\documentclass[Espanhol,Final]{tese-FT}

%Adicione seu arquivo com as referências bibliográficas
\addbibresource{bibliografia.bib}

%O pacote a seguir gera um dummy text. Elimine a linha quando
% for editar seu texto.
\usepackage{lipsum}

% Adicione o pacote a seguir, se seu texto tiver uma lista de 
% símbolos. Não confunda com a lista de abreviaturas.
%\usepackage[symbols,nogroupskip,sort=none]{glossaries-extra}

\begin{document}

% Escolha entre autor ou autora:
\autor{Nome do autor}
%\autora{Nome da Autora}

% Sempre deve haver um título em português:
\titulo{Título da Dissertação ou Tese em Português}

% Se a língua for o inglês ou o espanhol defina:
\title{The Dissertation or Thesis Title in English}

% Escolha entre orientador ou orientadora e inclua os títulos:
\orientador{Prof. Dr. Nome do Orientador}
%\orientadora{Profa. Dra. Nome da Orientadora}

% Escolha entre coorientador ou coorientadora, se houver, 
% e inclua os títulos:
%\coorientador{Prof. Dr. Eng. Lic. Nome do Co-Orientador}
\coorientadora{Prof. Dra. Eng. Lic. Nome da Co-Orientadora}

% Escolha entre uma das quatro opções a seguir (comente as demais):
%\bsi         % para Trabalho de Conclusão de Curso em BSI
%\tads       % para Trabalho de Conclusão de Curso em TADS
%\qualificacaoMestrado  % Para textos de qualificação de mestrado.
%\qualificacaoDoutorado % Para textos de qualificação de doutorado.
\mestrado   % para Dissertação de Mestrado em Tecnologia
%\doutorado  % para Tese de Doutorado em Tecnologia

%Defina a área de concentração. Se for TCC, deixe comentado
\areaConcentracao{Demografia}
%\areaConcentracao{Ambiente}
%\areaConcentracao{Ciência dos Materiais}

% Se houve cotutela, defina:
%\cotutela{Universidade Nova de Plutão}

%Defina a data da defesa no formato {Dia}{Mês}{Ano}
\datadadefesa{25}{02}{2019}

% Para a versão final defina:
% Repita o nome do Orientador(a) no primeiro avaliador
\avaliadorA{Prof. Dr. Nome do Orientador}{IFCH/UNICAMP}
\avaliadorB{Profa. Dra. Segunda Avaliadora}{Instituição da segunda avaliadora}
\avaliadorC{Dr. Terceiro Avaliador}{Instituição do terceiro avaliador}
% \avaliadorD{Prof. Dr. Quarto Avaliador}{Instituição do quarto avaliador}
% \avaliadorE{Prof. Dr. Quinto Avaliador}{Instituição do quinto avaliador}
% \avaliadorF{Prof. Dr. Sexto Avaliador}{Instituição do sexto avaliador}
% \avaliadorG{Prof. Dr. Sétimo Avaliador}{Instituição do sétimo avaliador}
% \avaliadorH{Prof. Dr. Oitavo Avaliador}{Instituição do oitavo avaliador}


% Para incluir a ficha catalográfica em PDF na versão final, 
% copie o arquivo PDF, descomente e ajuste a linha a seguir:
%\fichacatalografica{arquivo.pdf}

% Este comando deve ficar aqui:
\paginasiniciais

% Se houver dedicatória, descomente a linha a seguir:
%\prefacesection{Isso é um teste}
%A dedicatória deve ocupar uma única página.
%
%
% Se houver epígrafe, descomente e edite as linhas a seguir:
% \begin{epigrafe}
% {\it
% Vita brevis,\\
% ars longa,\\
% occasio praeceps,\\
% experimentum periculosum,\\
% iudicium difficile.}
%
% \hfill (Hippocrates)
% \end{epigrafe}
%
%
% Adicione no arquivo "agradecimentos.tex" os seus agradecimentos
% Caso prefira omitir os agradecimentos, comente a linha a seguir.
\input{agradecimentos}

% Sempre deve haver um resumo em português:

\input{resumo_port}

% Sempre deve haver um abstract:
\input{resumo_eng}

% Se houver um resumo em espanhol, descomente as linhas a seguir:
%\begin{resumen}
% A mesma regra aplica-se.
%\end{resumen}

% A lista de figuras:
\listoffigures

% A lista de tabelas:
\listoftables

% A lista de símbolos é opcional. Não confunda a lisa de símbolos
% com a lista de abreviaturas.
\input{listaSimbolos}

% A lista de abreviações e siglas vem a seguir.
% Dê uma olhada no pacote nomencl para ver os comandos para 
% adicionar abreviações e siglas no texto.
\renewcommand{\nomname}{Lista de Abreviações e Siglas}
\printnomenclature[3cm]

% O sumário vem aqui:
\tableofcontents

% E esta linha deve ficar bem aqui:
\fimdaspaginasiniciais

% O corpo da dissertação ou tese começa aqui:
%
% O comando a seguir inclui o arquivo introducao.tex
% que contém o capítulo de Introdução. 
% Detalhe: não precisa incluir a extensão .tex
\include{introducao}

% O comando a seguir inclui o arquivo revisao_bibliografica.tex
% que contém o capítulo de levantamento bibliográfico. 
% Detalhe: não precisa incluir a extensão .tex
\include{revisao_bibliografica}

% O comando a seguir inclui o arquivo metodologia.tex
% que contém o capítulo de metodologia. 
% Detalhe: não precisa incluir a extensão .tex
\include{metodologia}

% O comando a seguir inclui o arquivo resultados.tex
% que contém o capítulo de resultados. 
% Detalhe: não precisa incluir a extensão .tex
\include{resultados}

% O comando a seguir inclui o arquivo conclusoes.tex
% que contém o capítulo de conclusoes. 
% Detalhe: não precisa incluir a extensão .tex
\include{conclusoes}

% Os comandos para incluir as referências bibliográficas
\printbibliography[heading=bibintoc, % Adiciona no sumário
                   title=Referências bibliográficas % Nome do Capítulo
                  ]

% Os anexos, se houver, vêm depois das referências:
\appendix

% O comando a seguir inclui o arquivo apendices.tex
% que contém os apêndices. Observe o comando \appendix
% na linha anterior
% Detalhe: não precisa incluir a extensão .tex
\include{apendices}

\end{document}